\section{Algorithmes conçus}
\label{sec:algorithms}

Cette section présente, sous forme de pseudocode, les algorithmes tels
qu'ils sont réellement implémentés dans ce dépôt (\code{streamkm.core},
\code{streamkm.coreset}, \code{streamkm.spark}). Le code Scala correspondant,
avec ses commentaires, est présenté et discuté en
section~\ref{sec:implementation}.

\subsection{Seeding $D^2$ (k-means++) et algorithme de Lloyd}

Le seeding $D^2$ choisit les $k$ centres initiaux avec une probabilité
proportionnelle à la distance au carré au centre déjà choisi le plus proche,
garantissant une compétitivité $O(\log k)$ par rapport à l'optimum
\cite{arthur2007}. Notre implémentation (\code{KMeans.seedPlusPlus},
Algorithme~\ref{alg:seedpp}) en est une version \textbf{pondérée} : le
premier centre est tiré proportionnellement au poids des points plutôt
qu'uniformément (voir la discussion en section~\ref{sec:implementation}).

\begin{algorithm}[h]
\caption{\code{seedPlusPlus} : seeding $D^2$ pondéré}
\label{alg:seedpp}
\begin{algorithmic}[1]
\REQUIRE Points pondérés $\{(x_i, w_i)\}_{i=1}^n$, nombre de centres $k$
\ENSURE Centres initiaux $c_1,\dots,c_k$
\STATE $c_1 \leftarrow$ tirage parmi $x_i$ avec probabilité $\propto w_i$
\STATE $\mathrm{closest}[i] \leftarrow \lVert x_i - c_1\rVert^2$ pour tout $i$
\FOR{$f = 2$ à $k$}
  \STATE $c_f \leftarrow$ tirage parmi $x_i$ avec probabilité $\propto w_i \cdot \mathrm{closest}[i]$
  \FOR{tout $i$}
    \STATE $\mathrm{closest}[i] \leftarrow \min(\mathrm{closest}[i], \lVert x_i - c_f\rVert^2)$ \COMMENT{mise à jour incrémentale, évite $O(nk^2)$}
  \ENDFOR
\ENDFOR
\RETURN $(c_1,\dots,c_k)$
\end{algorithmic}
\end{algorithm}

L'algorithme de Lloyd affine ensuite ces centres par itérations
affectation/mise-à-jour jusqu'à convergence du coût. Notre implémentation
(\code{KMeans.lloyd}, Algorithme~\ref{alg:lloyd}) est également pondérée : le
centre d'un cluster est la \textbf{moyenne pondérée} de ses points affectés,
et non leur moyenne simple. Omettre les poids ici est le bug le plus
classique de ce pipeline, puisque le code continue de tourner sans erreur
mais produit des centres silencieusement faux dès qu'il est appliqué à un
coreset (poids hétérogènes).

\begin{algorithm}[h]
\caption{\code{lloyd} : algorithme de Lloyd pondéré}
\label{alg:lloyd}
\begin{algorithmic}[1]
\REQUIRE Points pondérés, centres initiaux $c_1,\dots,c_k$, tolérance $\tau$
\ENSURE Centres $c_1,\dots,c_k$, coût final
\STATE $\mathrm{prevCost} \leftarrow +\infty$
\REPEAT
  \STATE Affecter chaque point $x_i$ au centre le plus proche ; accumuler $\mathrm{currCost} \leftarrow \sum_i w_i \cdot \min_j \lVert x_i - c_j \rVert^2$
  \STATE Recalculer chaque centre $c_j$ comme la moyenne pondérée des points qui lui sont affectés
  \STATE \textbf{assert} $\mathrm{currCost} \le \mathrm{prevCost}$ \COMMENT{invariant de Lloyd, vérifié en mode debug}
  \IF{$\mathrm{prevCost} - \mathrm{currCost} \le \tau \cdot \mathrm{prevCost}$}
    \STATE \textbf{arrêter}
  \ENDIF
  \STATE $\mathrm{prevCost} \leftarrow \mathrm{currCost}$
\UNTIL{arrêt ou nombre maximal d'itérations atteint}
\RETURN centres, $\mathrm{currCost}$
\end{algorithmic}
\end{algorithm}

\code{KMeans.fit} exécute \code{nRestarts} exécutions indépendantes de
seeding + Lloyd et conserve le résultat de coût minimal, 5 par défaut,
conformément à la thèse (\cite{bitsakis2018}, section 2.3.4).

\subsection{Construction du coreset (coreset tree)}

\begin{algorithm}[h]
\caption{\code{CoresetTree.build} : construction du coreset par arbre binaire divisif}
\label{alg:coresettree}
\begin{algorithmic}[1]
\REQUIRE Points pondérés $P$, taille cible $m$
\ENSURE Coreset $S$ de taille $\le m$
\IF{$|P| \le m$}
  \RETURN $P$ \COMMENT{déjà son propre coreset}
\ENDIF
\STATE Racine $\leftarrow$ nœud contenant $P$, représentant tiré $\propto$ poids
\STATE $\mathrm{leaves} \leftarrow 1$
\WHILE{$\mathrm{leaves} < m$ et $\mathrm{cost}(\mathrm{racine}) > 0$}
  \STATE Descendre depuis la racine ; à chaque nœud interne, choisir un enfant avec probabilité $\propto$ son \textbf{coût}
  \STATE Dans la feuille atteinte, tirer un nouveau représentant $q$ avec probabilité $\propto D^2(\cdot, \mathrm{repr\_actuel})$
  \STATE Scinder la feuille en deux (proximité à l'ancien représentant vs. $q$)
  \STATE Recalculer coût et masse des deux nouveaux nœuds, propager la mise à jour jusqu'à la racine
  \STATE $\mathrm{leaves} \leftarrow \mathrm{leaves} + 1$
\ENDWHILE
\RETURN union des représentants des feuilles, pondérés par la masse de leur sous-arbre
\end{algorithmic}
\end{algorithm}

\subsection{Merge-and-reduce : consommation du flux}

\begin{algorithm}[h]
\caption{\code{BucketManager.insert} : insertion d'un point (merge-and-reduce)}
\label{alg:insert}
\begin{algorithmic}[1]
\REQUIRE Point $p$, taille de coreset $m$
\STATE $B_0 \leftarrow B_0 \cup \{p\}$
\IF{$|B_0| \ge m$}
  \IF{$B_1$ vide}
    \STATE $B_1 \leftarrow B_0$ ; $B_0 \leftarrow \emptyset$ \COMMENT{déplacement simple, pas de réduction}
  \ELSE
    \STATE $\mathrm{union} \leftarrow B_0 \cup B_1$ ; $B_0, B_1 \leftarrow \emptyset$ ; $\mathrm{level} \leftarrow 2$
    \WHILE{niveau $\mathrm{level}$ non trouvé libre}
      \STATE $\mathrm{reduced} \leftarrow \mathtt{CoresetTree.build}(\mathrm{union}, m)$
      \IF{$B_{\mathrm{level}}$ vide}
        \STATE $B_{\mathrm{level}} \leftarrow \mathrm{reduced}$ \COMMENT{niveau trouvé}
      \ELSE
        \STATE $\mathrm{union} \leftarrow \mathrm{reduced} \cup B_{\mathrm{level}}$ ; $B_{\mathrm{level}} \leftarrow \emptyset$ ; $\mathrm{level} \leftarrow \mathrm{level}+1$
      \ENDIF
    \ENDWHILE
  \ENDIF
\ENDIF
\end{algorithmic}
\end{algorithm}

\subsection{Portage Spark : fusion en arbre plutôt que séquentielle}

L'algorithme~\ref{alg:mergecoresets} est notre contribution principale par
rapport au dataflow Flink de la thèse : la fusion des coresets partiels de
toutes les partitions, qui y est un opérateur non-parallèle
(\code{FlatMapToFinalCoreset}, un seul thread), devient ici une réduction en
arbre distribuée.

\begin{algorithm}[h]
\caption{\code{StreamKMPlusPlus.mergeCoresets} : fusion distribuée en arbre}
\label{alg:mergecoresets}
\begin{algorithmic}[1]
\REQUIRE RDD de points $\mathrm{rdd}$, paramètres $(k,m,\mathrm{seed})$, profondeur $\mathrm{depth}$
\ENSURE Coreset final de taille $\le m$ représentant tout le RDD
\STATE Pour chaque partition d'indice $i$ : construire $\mathrm{bm}_i \leftarrow$ \texttt{BucketManager}$(m, \mathrm{seed}+i)$ et y insérer tous les points de la partition (Algorithme~\ref{alg:insert})
\STATE Fusionner les $\mathrm{bm}_i$ deux à deux en arbre, de profondeur $\mathrm{depth}$, via $\mathrm{bm}_a.\mathtt{merge}(\mathrm{bm}_b)$
\RETURN $\mathtt{extractCoreset}()$ du résultat de la fusion finale
\end{algorithmic}
\end{algorithm}

La graine dérivée de l'indice de partition (ligne 1) garantit un aléa
indépendant \emph{par partition}, une condition nécessaire pour que les
expériences de scalabilité (section~\ref{sec:experiments}) isolent l'effet
du partitionnement d'un simple effet de graine partagée (voir
section~\ref{sec:implementation} pour la discussion détaillée de ce choix
face à \code{RDD.treeAggregate}).
